\section{Scope and question addressed}

The present question is deliberately narrower than the full Heterogeneous Learning Systems
(HLS) programme. We ask only whether the following two theoretical components can be
supported individually and joined coherently.

\begin{description}[leftmargin=2.8cm,style=nextline]
\item[Beam 1.] Given a distribution of problems and currently available knowledge or
competence, how should work and access to expertise be organized?
\item[Beam 2.] How can performing work change the competence available in the future?
\end{description}

The intended feedback structure is
\begin{equation}
\text{current state}
\longrightarrow
\text{work organization}
\longrightarrow
\text{operational experience}
\longrightarrow
\text{future competence}
\longrightarrow
\text{future work organization}.
\label{eq:cycle}
\end{equation}

The goal of this note is \emph{not} to solve the dynamic optimization problem implicit in
\eqref{eq:cycle}. It is to establish whether the arrows in this skeleton have defensible
theoretical foundations and whether the variables at their interfaces can be given compatible
meanings.

Three distinctions are important from the outset.

First, a \emph{knowledge set}---the set of problems for which a solution is available---is
not automatically the same object as a scalar or continuous \emph{competence score}.
Second, the cost of acquiring knowledge is not competence itself. Third, an allocation
decision and the experience produced by executing that allocation are related but distinct
objects.

These distinctions prevent a superficial joining of otherwise different source theories.
